\subsection{General results}

Let X be a locally compact space in which a locally compact group H operates on the right, continuously and properly, by \((x,\xi)\mapsto x\xi\) ( \(x\in X\) , \(\xi\in H\) ). The equivalence relation in X defined by H is open (GT, III, §2, No.~4, Lemma 2) and X/H is Hausdorff (loc. cit., §4, No.~2, Prop. 3) hence locally compact (GT, I, §10, No.~4, Prop. 10). We denote by \(\pi\) the canonical mapping of X onto X/H. The saturation of a subset Y of X is \(\mathrm{YH}=\pi^{-1}\big(\pi(\mathrm{Y})\big)\) . If K is a compact subset of X, then \(\pi(\mathrm{K})\) is compact and the saturation \(\pi^{-1}\big(\pi(\mathrm{K})\big)\) of K is closed in X. Every compact subset of X/H is the image under \(\pi\) of a compact subset of X (GT, I, §10, No.~4, Prop. 10). We assume given once and for all a left Haar measure \(\beta\) on H.

Let \(\chi\) be a continuous representation of \(\mathbf{H}\) in \(\mathbf{R}_+^*\) . If a function \(g\) on \(\mathbf{X}\) satisfies \(g(x\xi) = \chi (\xi)g(x)\) for all \(x\in \mathbf{X}\) and \(\xi \in \mathbf{H}\) , its support \(\mathbf{S}\) is invariant under \(\mathbf{H}\) hence may be written \(\pi^{-1} (\pi (\mathrm{S}))\) . We shall denote by \(\mathcal{K}^{\chi}(\mathbf{X})\) the Riesz space formed by the continuous real-valued functions \(g\) on \(\mathbf{X}\) that satisfy \(g(x\xi) = \chi (\xi)g(x)\) ( \(x\in \mathbf{X},\xi \in \mathbf{H}\) ) and whose support is the saturation of a compact subset of \(\mathbf{X}\) ; we denote by \(\mathcal{K}_+^\chi (\mathbf{X})\) the set of elements \(\geqslant 0\) of \(\mathcal{K}^{\chi}(\mathbf{X})\) . In particular, \(\mathcal{K}^1 (\mathbf{X})\) is none other than the set of continuous functions on \(\mathbf{X}\) , constant on the orbits, whose support is the saturation of a compact subset.

PROPOSITION 1. --- Let f be a continuous real-valued function on X whose support S has compact intersection with the saturation of every compact subset of X.

\begin{enumerate}
\def\labelenumi{\alph{enumi})}
\tightlist
\item
  For every \(x \in \mathbf{X}\) , the function \(\xi \mapsto f(x\xi)\) on \(\mathbf{H}\) belongs to \(\mathcal{K}(\mathbf{H})\) ; one sets
\end{enumerate}

\[
f ^ {\chi} (x) = \int_ {\mathrm{H}} f (x \xi) \chi (\xi) ^ {- 1} d \beta (\xi).\tag{1}
\]

\begin{enumerate}
\def\labelenumi{\alph{enumi})}
\setcounter{enumi}{1}
\item
  The function \(f^{\chi}\) is continuous, is zero outside SH, and satisfies \(f^{\chi}(x\xi) = \chi(\xi)f^{\chi}(x)\) .
\item
  If \(g\) is a continuous real-valued function on \(\mathbf{X}\) and satisfies \(g(x\xi) = \chi(\xi)g(x)\) , then \((fg)^{\chi} = f^{1}g\) ( \(f^{1}\) being given by the formula (1) with \(\chi\) replaced by the representation \(\xi \mapsto 1\) of \(\mathbf{H}\) in \(\mathbf{R}_{+}^{*}\) ).
\item
  If \(\eta \in \mathbf{H}\) , then \(\left(\pmb{\delta}(\eta)f\right)^{\chi} = \chi (\eta)\Delta_{\mathrm{H}}(\eta)^{-1}f^{\chi}\) .
\end{enumerate}

Let \(x_0 \in X\) and let \(V\) be a compact neighborhood of \(x_0\) in \(X\) . The set of \(\xi \in H\) such that \(V\xi\) intersects \(S\) is also the set of \(\xi \in H\) such that \(V\xi\) intersects \(S \cap VH\) , hence is compact in \(H\) since \(S \cap VH\) is compact and \(H\) operates properly in \(X\) (GT, III, §4, No.~5, Th. 1); then Lemma 1 of §1, No.~1 proves a) and the continuity of \(f^{\chi}\) . The rest of b) is obvious. Finally, c) and d) result from the following calculations:

\[
\begin{array}{c} (f g) ^ {\chi} (x) = \int_ {\mathrm{H}} f (x \xi) g (x \xi) \chi (\xi) ^ {- 1} d \beta (\xi) = \int_ {\mathrm{H}} f (x \xi) g (x) \chi (\xi) \chi (\xi) ^ {- 1} d \beta (\xi) \\ = g (x) \int_ {\mathrm{H}} f (x \xi) d \beta (\xi) = g (x) f ^ {1} (x) \end{array}
\]

\[
\begin{array}{r l} & {\left(\delta (\eta) f\right) ^ {\chi} (x) = \int_ {\mathrm{H}} f (x \xi \eta) \chi (\xi) ^ {- 1} d \beta (\xi)} \\ & {\qquad = \Delta_ {\mathrm{H}} (\eta) ^ {- 1} \int_ {\mathrm{H}} f (x \xi) \chi (\xi \eta^ {- 1}) ^ {- 1} d \beta (\xi)} \\ & {\qquad = \chi (\eta) \Delta_ {\mathrm{H}} (\eta) ^ {- 1} \int_ {\mathrm{H}} f (x \xi) \chi (\xi) ^ {- 1} d \beta (\xi).} \end{array}
\]

PROPOSITION 2. --- The mapping \(f \mapsto f^{\chi}\) of \(\mathcal{K}(X)\) into \(\mathcal{K}^{\chi}(X)\) is linear, and the image of \(\mathcal{K}(X)\) (resp. \(\mathcal{K}_{+}(X)\) ) is \(\mathcal{K}^{\chi}(X)\) (resp. \(\mathcal{K}_{+}^{\chi}(X)\) ).

Linearity is immediate. It is clear that \(f^{\chi} \geqslant 0\) for \(f \geqslant 0\) . It then suffices to apply the following lemma:

Lemma 1. --- Let K be a compact subset of X, u a function of \(\mathcal{K}_{+}(X)\) with \(u(x) > 0\) for \(x \in K\) . Let \(g \in \mathcal{K}^{\chi}(X)\) be such that \(\operatorname{Supp} g \subset KH\) .

\begin{enumerate}
\def\labelenumi{\alph{enumi})}
\item
  One has \(\inf_{x\in \mathrm{KH}}u^1 (x) > 0\)
\item
  The function \(h\) equal to \(g / u^1\) on KH, and to 0 on X - KH, belongs to \(\mathcal{K}^{\chi}(X)\) .
\item
  \(g = (uh)^{\chi}\) .
\end{enumerate}

One has \(u^1(x) > 0\) for \(x \in \mathbf{K}\) , therefore \(\inf_{x \in \mathrm{KH}} u^1(x) = \inf_{x \in \mathbf{K}} u^1(x) > 0\) . Assertion b) follows from this at once. Finally, \((uh)^{\chi} = u^{1}h\) by Prop. 1 c), and it is clear that \(u^1h = g\) .

Let I be a relatively bounded linear form (Ch. II, §2, No.~2) on \(\mathcal{K}^{\chi}(\mathbf{X})\) . Then \(f \mapsto \mathrm{I}(f^{\chi})\) is a relatively bounded linear form on \(\mathcal{K}(\mathbf{X})\) , that is, a measure \(\mu_{\mathrm{I}}\) on \(\mathbf{X}\) . The mapping \(\mathrm{I} \mapsto \mu_{\mathrm{I}}\) is injective by Prop. 2. The measures \(\mu_{\mathrm{I}}\) on \(\mathbf{X}\) so obtained may be characterized as follows:

PROPOSITION 3. --- Let \(\mu\) be a measure on X. The following conditions are equivalent:

\begin{enumerate}
\def\labelenumi{\alph{enumi})}
\item
  There exists a relatively bounded linear form I on \(\mathcal{K}^{\chi}(\mathbf{X})\) such that \(\mathrm{I}(f^{\chi}) = \mu(f)\) for all \(f\in \mathcal{K}(\mathbf{X})\) .
\item
  \(\delta (\xi)\mu = \chi (\xi)^{-1}\Delta_{\mathrm{H}}(\xi)\mu\) for all \(\xi \in \mathbf{H}\) .
\item
  For all f, g in \(\mathcal{K}(\mathbf{X})\) ,
\end{enumerate}

\[
\mu (f \cdot g ^ {1}) = \mu (f ^ {\chi} \cdot g).\tag{2}
\]

\begin{enumerate}
\def\labelenumi{\alph{enumi})}
\setcounter{enumi}{3}
\item
  If \(f \in \mathcal{K}(\mathbf{X})\) is such that \(f^{\chi} = 0\) , then \(\mu(f) = 0\) .
\item
  \(\Rightarrow\) b): If \(\mu(f)=\mathrm{I}(f^{\chi})\) then, taking into account Prop. 1 d),
\end{enumerate}

\[
\begin{array}{r l} & {\langle \pmb {\delta} (\xi) \mu , f \rangle = \langle \mu , \pmb {\delta} (\xi^ {- 1}) f \rangle = \mathrm{I} \Big (\big (\pmb {\delta} (\xi^ {- 1}) f \big) ^ {\chi} \Big)} \\ & {\qquad = \mathrm{I} \big (\chi (\xi) ^ {- 1} \Delta_ {\mathrm{H}} (\xi) f ^ {\chi} \big)} \\ & {\qquad = \chi (\xi) ^ {- 1} \Delta_ {\mathrm{H}} (\xi) \langle \mu , f \rangle ,} \end{array}
\]

whence \(\delta (\xi)\mu = \chi (\xi)^{-1}\Delta_{\mathrm{H}}(\xi)\mu\)

\begin{enumerate}
\def\labelenumi{\alph{enumi})}
\setcounter{enumi}{1}
\tightlist
\item
  \(\Rightarrow\) c): Suppose hypothesis b) is satisfied. Note that the functions \((x,\xi)\mapsto f(x)g(x\xi)\) and \((x,\xi)\mapsto f(x\xi)g(x)\) on \(\mathbf{X}\times\mathbf{H}\) are continuous with compact support (because H operates properly in X); this established, Th. 2 of Ch. III, §4, No.~1 permits us to write:
\end{enumerate}

\[
\begin{array}{r l} & {\int_ {\mathrm{X}} f (x) d \mu (x) \int_ {\mathrm{H}} g (x \xi) d \beta (\xi) = \int_ {\mathrm{H}} d \beta (\xi) \int_ {\mathrm{X}} f (x) g (x \xi) d \mu (x)} \\ & {\qquad = \int_ {\mathrm{H}} d \beta (\xi) \int_ {\mathrm{X}} f (x \xi^ {- 1}) g (x) \chi (\xi) \Delta_ {\mathrm{H}} (\xi) ^ {- 1} d \mu (x)} \\ & {\qquad = \int_ {\mathrm{X}} g (x) d \mu (x) \int_ {\mathrm{H}} f (x \xi^ {- 1}) \chi (\xi) \Delta_ {\mathrm{H}} (\xi) ^ {- 1} d \beta (\xi)} \\ & {\qquad = \int_ {\mathrm{X}} g (x) d \mu (x) \int_ {\mathrm{H}} f (x \xi) \chi (\xi) ^ {- 1} d \beta (\xi),} \end{array}
\]

which proves c).

\begin{enumerate}
\def\labelenumi{\alph{enumi})}
\setcounter{enumi}{2}
\item
  \(\Rightarrow\) d): If c) is verified and if \(f^{\chi} = 0\) , then \(\mu(f \cdot g^{1}) = 0\) for all \(g \in \mathcal{K}(X)\) , thus \(\mu(f) = 0\) on choosing \(g \in \mathcal{K}(X)\) such that \(g^{1} = 1\) on \(\operatorname{Supp} f\) (which is possible by Prop. 2 applied with \(\chi = 1\) ).
\item
  \(\Rightarrow\) a): If condition d) is satisfied, there exists a linear form I on \(\mathcal{K}^{\chi}(\mathbf{X})\) such that \(\mu(f) = \mathrm{I}(f^{\chi})\) for \(f \in \mathcal{K}(\mathbf{X})\) , and this form is relatively bounded by virtue of Prop. 2.
\end{enumerate}

